\documentclass[11pt]{article}
\usepackage[T1]{fontenc}
\usepackage{lmodern}
\usepackage[margin=1.04in]{geometry}
\usepackage{microtype}
\usepackage{amsmath,amssymb,amsthm,mathtools,booktabs}
\usepackage[hidelinks]{hyperref}
\hypersetup{pdftitle={The fifth Taylor coefficient norm in Hardy spaces below H1}}
\newtheorem{theorem}{Theorem}
\newtheorem{proposition}[theorem]{Proposition}
\newtheorem{lemma}[theorem]{Lemma}
\newcommand{\D}{\mathbb D}
\newcommand{\T}{\mathbb T}
\newcommand{\Dif}{\mathcal D}
\newcommand{\coef}[1]{[z^{#1}]}
\allowdisplaybreaks[1]
\title{The fifth Taylor coefficient norm\\in Hardy spaces below \texorpdfstring{\(H^1\)}{H1}}
\author{}
\date{27 November 2026}
\begin{document}
\maketitle

\begin{abstract}
We prove a uniform single-root formula for the sharp fifth Taylor
coefficient norm on \(H^p(\mathbb D)\) for \(0<p<1\).
Every extremizer has a positive real polynomial factor after scalar and
argument rotations. The formula uses a univariate polynomial of
degree fifteen and two polynomial reconstruction equations. The
maximizing factor is unique up to rotations and depends real
analytically on \(p\). The phase proof uses a
second-order variation when \(2/3<p<1\) and a reciprocal-root
bound when \(1/2<p<2/3\). The underlying algebraic reduction
bounds the number of nonreal stationary pairs by five, even
before the second-order test. We also evaluate the endpoint slope
at \(p=1\) and prove uniqueness of the analytic maximizing branch
near that endpoint. The sharp constants have algebraic degrees
60 and 52 at \(p=4/7\) and \(p=4/5\), respectively. Neither is
expressible by radicals.
\end{abstract}

\section{The fifth coefficient problem}

Let \(m\) be normalized arc measure. For \(p>0\), put
\[
 \|f\|_p=\sup_{0<r<1}
       \bigg(\int_{\T}|f(rz)|^p\,dm(z)\bigg)^{1/p},\qquad
 C(n,p)=\sup_{\|f\|_p\leq1}|[z^n]f|,
\]
where \(f\) is holomorphic on \(\D\). Brevig and Saksman
\cite{BrevigSaksman} formulated the sharp coefficient problem for
\(0<p<1\). They determined \(C(2,p)\) and \(C(3,2/3)\).
The outer theorem of Brevig and S{\o}vik \cite{BrevigSovik}
applies to this problem in every degree.

Earlier work on the present problem computed the third and
fourth coefficients, and then found a finite fifth-coefficient
formula for \(0<p\leq1/2\) and isolated the value at \(p=2/3\)
\cite{Low,High}. We use the same reflected stationary identity
but allow arbitrary phases at the beginning. The auxiliary
problem is to bound the number of nonreal stationary polynomials
in a fixed exponent. The third and fourth Taylor identities
determine all nonreal coefficients from one real modulus, and
the fifth identity gives a quintic. A second variation then
turns any nonreal maximizing point at \(2/3<p<1\) into a
continuum of such points. On \(1/2<p<2/3\) the reversed
polynomial gives a different exclusion. This leaves a fixed
univariate elimination for the real candidates in all exponents.
The approach draws on the jet restrictions and phase analysis
in the Verifiable Mathematical Reasoning Corpus \cite{VMRC}.

\section{Boundary separation and stationary jets}

\begin{proposition}\label{prop:boundary}
Fix \(n\ge1\), \(0<p<1\), and \(s=2/p\). The norm \(C(n,p)\)
is attained. Every extremizer of quasi-norm one has the form
\[
 f=\eta Q^s,\qquad
 \deg Q=n,\quad \|Q\|_2=1,\quad Q(0)>0,\quad
 Q(z)\ne0\ \text{for }|z|\le1,\quad |\eta|=1 .
\]
The power uses the analytic logarithm of \(Q\) which is real at zero.
\end{proposition}

\begin{proof}
The Hardy pointwise estimates give a normal family on every smaller
disk. A locally uniform limit of a maximizing sequence still has
quasi-norm at most one, by integration on each smaller circle.
The coefficient converges and the maximum is positive, since \(z^n\)
is admissible. Scaling shows every extremizer has quasi-norm one.
Brevig--S{\o}vik's outer theorem \cite[Thm.~1]{BrevigSovik} applies
to the bounded linear functional \(\coef n\). For convenience, its
concavity argument is as follows. Write \(f=IO\), where \(I\) is
inner and \(O\) is outer. If \(I\) is nonconstant, set
\(u=2\Re I-c\) on \(\T\), with
\(c=\int_{\T}2\Re I\,|f|^p\,dm\). Then \(u\) is nonconstant and has
mean zero under \(|f|^p\,dm\). The analytic function
\(g=O(1+I^2-cI)\) belongs to \(H^p\), and \(g=uf\) on \(\T\).
For small real \(t>0\), strict concavity of \(x^p\) gives
\(\|f\pm tg\|_p^p<1\), since \(1\pm tu>0\).
Averaging the two coefficient values contradicts extremality.
Thus \(f\) is outer. Here \(\Re I\) cannot be constant on \(\T\)
unless \(I\) is constant, and the outer boundary modulus is positive
almost everywhere.

By rotations assume \(a_0=f(0)>0\) and \(a_n=\coef n f=C(n,p)>0\).
Vary \(f e^{t g}/\|f e^{t g}\|_p\), first with an analytic polynomial
\(g\) and then with \(ig\). Differentiating the norm is justified
by boundedness of \(g\) on \(\T\). Stationarity gives
\[
 \coef n(fg)=a_n\int_{\T}|f|^p g\,dm .
\]
With \(g=z^k\) for every \(k\ge0\), Fourier uniqueness gives
\[
 h(z)=a_n+\sum_{k=1}^n
       \bigl(a_{n-k}z^{-k}+\bar a_{n-k}z^k\bigr)
     =a_n |f(z)|^p \quad\text{on }\T\text{ a.e.}
\]
The Fej\'er--Riesz theorem factors \(h=a_n |Q|^2\) with
\(Q\) free of zeros in \(\D\), \(Q(0)>0\), \(\|Q\|_2=1\), and
degree exactly \(n\) because \(a_0>0\). The normalized outer function
with boundary modulus \(|Q|^{2/p}\) is \(Q^s\), so \(f=Q^s\).

If \(U\) has degree at most \(d\), write \(U^{*d}(z)=
z^d\overline{U(1/\bar z)}\).
Let \(R=z^n h=a_n QQ^{*n}\), let \(\Dif=z\,d/dz\), and set
\(E=R\Dif\log f=s a_n Q^{*n}\Dif Q\).
The polynomial \(R\) is self-reciprocal of degree \(2n\), and
\(R=f\bmod z^{n+1}\), meaning their Taylor coefficients agree through
order \(n\). Consequently the first \(n+1\) coefficients of \(E\)
are \(k a_k\), \(0\le k\le n\). The product rule and
\((\Dif Q)^{*n}=nQ^{*n}-\Dif(Q^{*n})\) give
\[
 E-E^{*2n}=s(\Dif R-nR).
\]
Set \(\beta=1-1/p=1-s/2\). The first \(n+1\) coefficients of
\((E+E^{*2n})/2\) are \((n-\beta(n-k))a_k\).
Reflecting the remaining coefficients shows that, on \(\T\),
\begin{equation}\label{eq:boundary-identity}
 z^{-n}\frac{E+E^{*2n}}{2}=nh-\beta D_h,
\end{equation}
where \(D_h\) is the outward radial derivative of the harmonic
extension of \(h\), the Fourier multiplier \(|k|\).

Suppose \(Q(\zeta)=0\) for some \(\zeta\in\T\). Then \(E(\zeta)=0\)
by its factored form and \(E^{*2n}(\zeta)=0\) by reflection.
After rotation to \(\zeta=1\), write \(Q(z)=(z-1)B(z)\).
Applying the multiplier to
\(h=a_n(2-z-z^{-1})|B|^2\) at 1 gives
\(D_h(1)=-2a_n\|B\|_2^2<0\), since
\(2|k|-|k+1|-|k-1|\) vanishes except at \(k=0\), where it
equals \(-2\). Equation \eqref{eq:boundary-identity} would now give
\(0=-\beta D_h(\zeta)\), which is impossible. This proves strict
boundary separation.
\end{proof}

There are two useful forms of stationarity. Put \(a=s-1>1\),
\(q_0=Q(0)\), and \(P=Q/q_0\). With \(a_n>0\) as in the proof,
the leading coefficient \(d=\coef n P\) is positive because the
constant term of \(R=a_n QQ^{*n}\) is \(q_0^s>0\).
The congruence \(R=f\bmod z^{n+1}\) gives
\begin{equation}\label{eq:jet}
 P^a\equiv P^{*n}/d\pmod{z^{n+1}},\qquad
 a_n=\frac{q_0^{a-1}}{d},\qquad
 q_0=(1+|[z]P|^2+\cdots+|\coef nP|^2)^{-1/2}.
\end{equation}
When rotating the argument without the \(a_n>0\) normalization,
\(P^a\equiv P^{*n}/\bar d\bmod z^{n+1}\) and
\(\coef n Q^s=q_0^{a-1}/\bar d\). These identities also follow
from \eqref{eq:jet} by rotation.
The boundary argument also applies to a formal solution of the
jet with \(P(0)=1\) and \(d\ne0\). To see this, multiply the
jet by \(P\), take the formal power \(f=P^{a+1}\) at zero, and
use \(R=PP^{*n}/\bar d\) after choosing \(d>0\) by rotation.
The proof of \eqref{eq:boundary-identity} then uses only Taylor
coefficients and the finite trigonometric polynomial
\(z^{-n}R=|P|^2/d\) on \(\T\). It therefore still excludes a
unit-circle root. This observation controls limits of solutions
when the exponent varies.

\section{A degree-fifteen formula}

Use \(a=2/p-1\), \(B=a-2\), and \(\delta=a-1=B+1>0\).
For \(w\in\mathbb R\), define
\[
 \begin{gathered} k=B/6,\qquad Q(w)=w^2+Bw+B(B-1)/12, \\
 R(w)=Bw^2/2+B(B-1)w/6+B(B-1)(B-2)/120.\end{gathered}
\]
Here \(Q(w),R(w)\) denote scalar polynomials in \(w\).
Set \(L=y(w+\delta/2)+\delta(w+k)\) and
\begin{align}
 H_4(w,y)={}&2yL^2-2y^2wL-\delta w^2\bigl(2y(w+\delta/2)+Q(w)\bigr),
       \label{eq:real4}\\
 H_5(w,y)={}&a y^2L^3-y^3wL^2-a\delta y^2w^2L \notag\\
     &\quad-a\delta w^3\bigl(y(w+\delta/2)(w+B/2)+R(w)\bigr).
       \label{eq:real5}
\end{align}
Define the univariate polynomial
\begin{equation}\label{eq:real-elim}
 \mathcal R_B(w)=\frac{\operatorname{Res}_y(H_4,H_5)}{w^9(w+\delta/2)^5}.
\end{equation}
The resultant in \eqref{eq:real-elim} is the determinant of the
Sylvester matrix with \(H_4\) first. Its degree in \(w\) is
twenty-nine and the denominator divides it identically.
Consequently \(\mathcal R_B\) has degree fifteen and leading
coefficient \(-\delta^8\).

\begin{theorem}\label{thm:main}
For \(0<p<1\), form the following list.

For each positive root \(w\) of \(\mathcal R_B\), take
each \(y>0\) satisfying \(H_4(w,y)=H_5(w,y)=0\) and
\(L>0\). Set \(X=w/(yL)\), \(x=\sqrt X\), \(d=xy/a\), and
\begin{equation}\label{eq:real-P}
 P(z)=1+xz+Xwz^2+x^3y(w+\delta/2)z^3+Xyz^4+dz^5.
\end{equation}
Keep \(P\) if all zeros of \(P\) lie outside \(\overline{\D}\).
The list is nonempty, contains at most forty-five polynomials, and
\[
 \boxed{\displaystyle C(5,p)=
     \max_{\text{\(P\) kept}}\frac{\|P\|_2^{-(a-1)}}{d}.}
\]
All unit-quasi-norm extremizers are unimodular multiples and
argument rotations of \((P/\|P\|_2)^{2/p}\) for maximizing entries.
\end{theorem}

The derivation below proves that every maximizing factor can be
rotated to have strictly positive coefficients. Once this is known, the
resultant separates the remaining possibilities without tracking
one branch as \(p\) varies.

\begin{proof}[Proof of Theorem~\ref{thm:main} given phase positivity]
Write \(P=1+xz+bz^2+cz^3+ez^4+dz^5\).
If \(x=0\), the first jet gives \(e=0\), the fourth gives
\(\binom{a}{2}b^2=0\), and the second gives \(c=0\).
Then \(P^{2/p}\) is a function of \(z^5\), and its fifth
coefficient cannot exceed \(C(1,p)\) after norm normalization.
However \(C(1,p)<C(5,p)\): multiply a first-coefficient extremizer
by \(z^4\) and use the outer theorem to rule out equality.
We may therefore rotate to \(x>0\). The positivity proof below shows that \(w,y>0\).

Define real \(w=b/x^2\) and \(y=ad/x\), and put \(X=x^2\).
The first two jets give \(e=Xy\) and \(c=x^3y(w+\delta/2)\).
With \(c,e,d\) substituted, the remaining three identities are
\begin{align*}
 w&=XyL,\\
 1&=y^2 X+\delta X^2\bigl(y^2(w+\delta/2)+yQ(w)/2\bigr),\\
 1&=\frac{y^2X}{a}+\delta y^2X^2+
        \delta X^3\bigl(y^2(w+\delta/2)(w+B/2)+yR(w)\bigr).
\end{align*}
These follow by expanding \(P^a\) through order five, and are
equivalent to the remaining jets. The first identity gives \(X=w/(yL)\). Substituting it in the
other two and clearing the nonzero denominators gives
\eqref{eq:real4} and \eqref{eq:real5}. At \(w,y>0\), the
division in \eqref{eq:real-elim} removes only nonzero factors.

For each positive \(w\), \(H_4\) is a cubic in \(y\) with
leading coefficient \(\delta(w+\delta/2)>0\). Thus
\(\mathcal R_B\) gives at most fifteen positive \(w\), each
with at most three positive \(y\).

With nonzero denominators and \(X>0\) all steps are reversible,
so each retained \(P\) satisfies the stationary jet.
Once \(P\) is zero-free on \(\overline{\D}\), its power is
an admissible comparison function of unit quasi-norm after
normalization. Multiplying the jet by \(P\) gives its coefficient
modulus as \(\|P\|_2^{-(a-1)}/|d|\), using
\([z^5]PP^{*5}=\|P\|_2^2\). Every maximizer occurs in the list
and every entry is admissible, which proves the formula and
the equality cases.
\end{proof}

\subsection{Exact root tests}

The disk-root test in Theorem~\ref{thm:main} can be carried out
by five polynomial steps. Start with the monic reciprocal
polynomial \(F=P^{*5}\). For a monic polynomial \(F\) of
degree \(j\), let \(c=F(0)\). If \(|c|\geq1\), reject it.
Otherwise replace \(F\) by the monic polynomial
\[
 \frac{F-cF^{*j}}{z(1-|c|^2)} .
\]
Repeat until degree zero. If \(F\) has all roots in \(\D\),
then \(|c|<1\) and Rouch\'e's theorem gives the corresponding
assertion for the replacement. Conversely a root of \(F\) on
\(\T\) would be a root of its replacement. When the replacement
has all roots in \(\D\), \(F\) has no root on \(\T\), so
Rouch\'e's theorem again proves the converse. This yields the
claimed exact test with strict inequalities at every step.

For rational \(p\), isolating the real roots of the two displayed
polynomials and comparing their algebraic values gives a terminating
calculation. The same is true for algebraic \(p\). In the irrational
algebraic case a comparison of distinct normalization factors
reduces to comparing a nonzero algebraic number with a positive
algebraic base other than \(1\), raised to the irrational algebraic
power \((a-1)/2\). Gelfond--Schneider excludes equality, and certified
approximations then decide the strict comparison. Equal normalization
factors reduce directly to algebraic comparison. The formula itself
holds for every real \(p\) in the interval.

\section{Elimination of nonpositive phases}

The proof uses the parameter \(a=2/p-1\). We treat the two ranges in
which the higher binomial coefficients change sign, along with the
integer anchors and the range \(a>3\).

\subsection{A quintic bound for the nonreal candidates}
\label{sec:nonreal-candidates}
Fix \(1/2<p<1\) and write \(a=2/p-1\), \(B=a-2\),
\(\delta=B+1\). The following computation classifies the stationary
polynomials whose \(z^2\) coefficient is nonreal after rotation to
make the \(z\) coefficient positive. In particular, each exponent has
at most five conjugate pairs of these normalized candidates.

Use the gauge \(x>0\) just established for maximizing factors. Normalize the constant coefficient to one and
rotate so that the first coefficient \(x\) is positive. Put
\(X=x^2\) and write the degree-five polynomial as
\[
 P=1+xz+x^2\xi z^2+x^3y(\bar\xi+\delta/2)z^3+
           x^2yz^4+(xy/a)z^5 .
\]
The first two jets give the displayed \(z^3,z^4\) coefficients.
Define
\[
 \begin{gathered}k=B/6,\quad Q(\xi)=\xi^2+B\xi+B(B-1)/12,\\
 R(\xi)=B\xi^2/2+B(B-1)\xi/6+B(B-1)(B-2)/120,\quad r=\frac1{X|y|^2}-1.\end{gathered}
\]
Here \(y\ne0\) and \(r>-1\). The third through fifth jets read
\begin{align*}
 \delta(\xi+k)/y&=r\bar\xi-\delta/2,\\
 \frac{r}{X\delta}&=\bar\xi+\delta/2+Q(\xi)/(2y),\\
 1+r&=\frac1a+\delta X+
    \delta X^2(\bar\xi+\delta/2)(\xi+B/2)+\delta X^2R(\xi)/y.
\end{align*}
The last jet, using
\(R=(B-2)Q/10+(2B+1)\xi(\xi+k)/5\) and the preceding jets, is
the single real equation
\begin{equation}\label{eq:c5-missing-fifth-real}
 0=\frac1a+X\!\left(\delta+\frac{(B-2)r}{5}\right)
  +X^2\left[\left(\delta+\frac{(2B+1)r}{5}\right)|\xi|^2+
       \frac{\delta(3B+4)}5\Re\xi+\frac{\delta^2(3B+4)}{20}\right]
  -1-r .
\end{equation}
In the nonreal case \(\xi+k\ne0\). Put \(t=\Re\xi+k\),
\(K=|\xi+k|^2\), \(m=rk+\delta/2\), and \(\vartheta=(2B+3)/30\).
Use the third jet to substitute for \(1/y\) in the fourth.
Taking imaginary parts and dividing by \(\Im\xi\) gives
\[
 B\left[-r-\frac{\vartheta(m-2rt)}K\right]=3\delta .
\]
One way to check this calculation is by the division
\[
 \frac{Q(\xi)}{\xi+k}=\xi+\frac{5B}{6}
                        -\frac{5B\vartheta}{6(\xi+k)}.
\]
Thus \(B=0\) admits no nonreal candidates. Otherwise define
\(h=(r+3\delta/B)/\vartheta\). We have \(m-2rt=-hK\) and
\[
 |r\bar\xi-\delta/2|^2=r^2K-2rmt+m^2=(r^2-mh)K.
\]
Since \(X|y|^2=(1+r)^{-1}\), taking norms in the third jet gives
\[
 W=r^2-mh>0,\qquad X=\frac{W}{\delta^2(1+r)}.
\]
The real part of the fourth jet is
\[
 \frac{2r}{X}-\delta^2=rK-2mt+\frac{B(3-2B)}{36}r-\frac{2B\delta}3.
\]
Indeed, \(\Re[(r\bar\xi-\delta/2)/(\xi+k)]=ht-r\);
substitution in the displayed division gives the identity.
If \(r=0\), the imaginary equation would give \(K=-B\vartheta/6\),
requiring \(B<0\), while the real equation gives \(t=1+B/3\).
These violate \(K>t^2\). Hence \(r\ne0\). We obtain
\begin{equation}\label{eq:c5-nonreal-reconstruct}
 K=\frac{2r^2/X-\delta^2r+m^2-B(3-2B)r^2/36+2B\delta r/3}{W},
 \qquad t=\frac{m+hK}{2r}.
\end{equation}

For each \(B\), substitute \(|\xi|^2=K-2kt+k^2\) and
\(\Re\xi=t-k\) in \eqref{eq:c5-missing-fifth-real}, and replace
\(X,K,t\) by their displayed rational formulas. Clearing
denominators yields a polynomial \(\sum_{j=0}^5 c_jr^j\),
where
\begin{align*}
 c_5={}&3B^3(B-1)^3(B+2)(2B+1),\\
 c_4={}&-60B^3(B-1)(B+1)(B+2)(B^2+7B+2),\\
 c_3={}&-10B^2(B+1)^2(B+2)(8B^4+70B^3+77B^2-456B-99),\\
 c_2={}&-20B(B+1)^3(12B^5+50B^4-307B^3-1929B^2-2556B-270),\\
 c_1={}&-5B(B+1)^4(48B^4-422B^3-6045B^2-19611B-19170),\\
 c_0={}&-20(B+1)^5(4B^4-213B^3-1791B^2-4725B-4050).
\end{align*}
More precisely its quotient by
\(20B^2(B+1)^4(B+2)(2B+3)^2(1+r)^2\) equals the left side of
\eqref{eq:c5-missing-fifth-real} after the substitution.
The leading coefficient \(c_5\) is nonzero except possibly
at \(B=-1/2\), where \(c_4\ne0\). Thus there are at most
five real roots to check.

For each such root retain only \(r>-1\), \(r\ne0\), \(W>0\)
and \(K>t^2\). Reconstruct
\[
 \xi=t-k+i\sqrt{K-t^2},\qquad
 y=\frac{\delta(\xi+k)}{r\bar\xi-\delta/2},\qquad x=\sqrt X.
\]
The denominator in \(y\) is nonzero because its squared modulus
is \(WK>0\). The conjugate reconstruction gives the only other
candidate for that root. The formulas are reversible and imply
all five stationary jets. Retain \(P\) when all of its zeros
lie outside the closed disk. Its admissible coefficient value is
\(\|P\|_2^{-(a-1)}/|xy/a|\).


\subsection{Excluding nonreal extremizers for \texorpdfstring{\(1<a<2\)}{1<a<2}}
\label{sec:low-hessian}

We first eliminate a vanishing phase term. Normalize \(P(0)=1\)
and rotate \(x=[z]P\) to be positive. The first two jets give
\(e=axd,\ c=d(a\overline b+\alpha x^2)\), where
\[
 \alpha=\binom a2,\quad \beta=\binom a3,\quad
 \gamma=\tfrac15\binom a4,\quad
 Y=x(\alpha b^2+\beta x^2b+\gamma x^4).
\]
Up to the factor \(s q_0^s\), the fifth coefficient is
\(d+axe+(ab+\alpha x^2)c+Y\). When \(a\ne2\), all zeros of
\(Y\), considered as a polynomial in \(b\), are simple: its
discriminant apart from \(x^6\) is
\(\beta^2-4\alpha\gamma=\alpha^2(a-2)(2a-1)/45\).
Suppose \(Y=0\) at a maximizer. Rotate \(d,e,c\) by a common
phase \(\phi\), and independently vary \(b\) along a complex
direction \(h\) with real parameter \(\tau\). Use the \(H^2\)
normalizing factor for each perturbation. On the base slice
\(\tau=0\), the squared objective is constant in \(\phi\).
At the base point the mixed \(\phi,\tau\) derivative is a
nonzero real factor times
\(\Re(iZ\overline{Y'(b)h})\), where
\(Z=d+axe+(ab+\alpha x^2)c=dH\ne0\) and
\(H=1+a^2x^2+|ab+\alpha x^2|^2\).
Choose \(h\) to make this real part nonzero. An admissible
local maximum would have a negative semidefinite Hessian,
which cannot have a zero diagonal entry and a nonzero entry
in its corresponding row. Thus \(Y\ne0\).
Rotating the three high coefficients together now shows that
the phases of \(d\) and \(Y\) agree.

Put \(t=x^2,\ w=b/t,\ U=|w|^2,\ W=\Re w\), and fix
\(T=\|P\|_2^2-1-t\) and \(t\). If \(w\) is nonreal, \((U,W)\)
are local coordinates and \(U>W^2\). Write
\[
 H_0=a^2U+\alpha^2+2a\alpha W,\quad H=1+a^2t+t^2H_0,
 \qquad f=|\alpha w^2+\beta w+\gamma|^2 .
\]
For nearby \(U,W\), realize the fifth-coefficient modulus,
apart from the unchanged factor \(s q_0^s\), as
\begin{equation}\label{eq:fifth-slice-geometry}
 \mathcal F(U,W)=\sqrt{(T-t^2U)H}+t^{5/2}\sqrt f .
\end{equation}
Indeed, choose \(w=W+i\,\operatorname{sgn}(\Im w_{\rm base})
\sqrt{U-W^2}\), choose
\(d=\sqrt{(T-t^2U)/H}\;Y/|Y|\), and set
\(e=axd,\ c=d(a\overline b+\alpha x^2)\).
These give the original factor at the base point and are
admissible nearby by strict boundary exclusion.

In the range \(1<a<2\), the second term of
\eqref{eq:fifth-slice-geometry} is a positive constant times
the norm from a positive definite quadratic form centered at
\[
 (U_0,W_0)=\left(\gamma/\alpha,-\beta/(2\alpha)\right).
\]
In fact
\[
 f=\alpha^2(U-U_0)^2+
      2\alpha\beta(U-U_0)(W-W_0)+4\alpha\gamma(W-W_0)^2,
 \qquad 4\alpha\gamma-\beta^2>0.
\]
Away from the center its Hessian is positive semidefinite of
rank one, with radial null direction. The first term is the
geometric mean of two positive affine functions of \(U,W\).
Its second derivative in direction \(\xi\) equals
\[
 -\frac{\sqrt{\ell_1\ell_2}}4
  \left(\frac{\ell_1'(\xi)}{\ell_1}
                  -\frac{\ell_2'(\xi)}{\ell_2}\right)^2,
 \quad \ell_1=T-t^2U,\quad \ell_2=H .
\]
The normalized gradients are distinct since
\(\ell_1\) is independent of \(W\) and \(\partial_W H>0\).
Thus its Hessian is negative semidefinite of rank one.
At a local maximum the two null directions must coincide, since
otherwise the second derivative in the first term's null
direction would be positive. Along the common null line both
summands are affine near the base point. For the first this
follows from proportional relative variations of its two
affine factors. For the second it follows because the line
passes through \((U_0,W_0)\) and the base point is distinct from
the center. The derivative of their sum vanishes at the base
point, so \(\mathcal F\) is constant along a small interval in
this line. Our continuous realization then gives infinitely
many distinct nonreal maximizers in the \((U,W)\) coordinates.
This contradicts the finite enumeration of nonreal stationary
candidates proved above. Consequently \(\xi=b/x^2\) is real at every maximizer in this
range. In the notation of Section~\ref{sec:nonreal-candidates},
the third jet then makes \(y\) real unless \(\xi=-k\).
In that exceptional case \(Q(-k)=-B(2B+3)/36\ne0\), and the
fourth jet makes \(1/y\) real. The first two jets give real
\(c,e\) as well.


\subsection{Positivity of the real factor for \texorpdfstring{\(1<a<2\)}{1<a<2}}

Assume now that the normalized maximizing factor is real. We show
that every coefficient is positive. Write
\[
 \eta=2-a,\quad \delta=a-1=1-\eta,\quad
 B=\eta/3,\quad G=\eta(1+\eta)/60,\quad v=ad,\quad X=x^2.
\]
Here \(Y=\alpha x^5(w^2-Bw+G)\) with \(w=b/x^2\) and
\(4G>B^2\). Hence \(Y>0\) when \(w\) is real, so phase
alignment gives \(d>0\).
Suppose \(w=-u\) for \(u\geq0\). If \(u>0\), allow the phase
of \(w\) to vary, with \(U=u^2\) fixed in
\eqref{eq:fifth-slice-geometry}. The maximum at \(W=-u\)
implies that the one-sided derivative with respect to \(W\)
is nonpositive. Evaluating this derivative yields
\[
 0\geq \alpha X^2\left(
 v-x\,\frac{B(u^2+G)+4Gu}{u^2+Bu+G}\right).
\]
Indeed the derivative of the first term is \(\alpha X^2 v\).
The derivative of the second follows by differentiating
\(f=(\alpha U-\gamma)^2+\beta^2U+
 2\beta(\alpha U+\gamma)W+4\alpha\gamma W^2\).
Since \(u^2+G\geq2u\sqrt G\), we have
\[
 \frac{B(u^2+G)+4Gu}{u^2+Bu+G}
 =B+\frac{(4G-B^2)u}{u^2+Bu+G}\leq 2\sqrt G.
\]
If \(u=0\), the third jet instead gives \(v/x=B\).
In each case \(k=v/x\leq2\sqrt G\).

Strict boundary exclusion gives \(X(1+2u)<5\). In fact, if
\(P(z)=\prod_{i=1}^5(1-\lambda_i z)\) with \(|\lambda_i|<1\),
then \(\sum_i\lambda_i^2=x^2-2b=X(1+2u)\).
On the other hand, substitute \(b=-uX\), \(v=kx\),
\(e=vx\), and \(c=v(b+\delta X/2)\) into the fifth jet.
After multiplying by \(d=v/a\) it is
\[
 1=\frac{k^2X}{a}+\delta k^2X^2+
 \delta X^3\left[
 k^2(\delta/2-u)(-u-\eta/2)
 -k\left(\frac{\eta u^2}{2}
        +\frac{\eta(1+\eta)u}{6}
        +\frac{\eta(1+\eta)(2+\eta)}{120}\right)\right].
\]
Dropping the nonpositive summand and writing \(r_+=\max(r,0)\)
gives
\[
 1\leq k^2\left[\frac{X}{a}+\delta X^2+
         \delta X^3(u-\delta/2)_+(u+\eta/2)\right].
\]
Use \((u-\delta/2)_+(u+\eta/2)\leq u(1+2u)/2\) and
\(X<5/(1+2u)\). The bracket is bounded above by
\[
 \frac5a+25\delta\,\frac{1+5u/2}{(1+2u)^2}
 \leq \frac5a+25\delta .
\]
The resulting upper bound on the right side is
\[
 4G(5/a+25\delta)
 =\frac{\eta(1+\eta)}{15}
      \left(\frac5{2-\eta}+25(1-\eta)\right)<1.
\]
For the last inequality, multiplication by \(3(2-\eta)\)
reduces it to the positivity of
\(6-14\eta+4\eta^2+10\eta^3-5\eta^4\).
In the degree-four Bernstein basis on \([0,1/2]\) its coefficients
are \(6,17/4,8/3,25/16,15/16\), and on \([1/2,1]\) they are
\(15/16,5/16,1/6,1/2,1\). All are positive. The contradiction
proves \(b>0\), and \(c=v(b+\delta X/2)>0\) and \(e=vx>0\)
complete the claim.


\subsection{The reciprocal-root exclusion for \texorpdfstring{\(2<a<3\)}{2<a<3}}

\begin{lemma}\label{lem:middle-real-phase}
Suppose \(2<a<3\). If
\(P(z)=1+xz+bz^2+cz^3+ez^4+dz^5\) is zero-free on the closed
unit disk and satisfies the reflected stationary jet
\(P^a\equiv P^{*5}/\bar d\pmod {z^6}\), with \(x>0\), then
\(b,c,e,d\) are real.
\end{lemma}

\begin{proof}
In this proof put \(\epsilon=a-2\), \(\delta=a-1\), \(w=b/x^2\),
\(u=\Re w\), \(W=|w|^2\), \(X=x^2\), \(V=a^2|d|^2\),
\(\ell=(1-V)/V\), and \(h=a\bar d x/V\).
The first two jets give \(e=axd\) and
\(c=d(a\bar b+\binom a2 x^2)\). The next two give
\begin{align}
 \ell\bar w&=\delta/2+\delta h(w+\epsilon/6), \label{eq:mid-q-3}\\
 \ell&=\delta X\left(\bar w+\delta/2+h B\right),
 \qquad B=w^2/2+\epsilon w/2+\epsilon(\epsilon-1)/24 .
 \label{eq:mid-q-4}
\end{align}
If \(w\) is nonreal, solve \eqref{eq:mid-q-3} for \(h\) and set
\begin{align*}
 D&=60W+12(\epsilon-1)u+\epsilon(\epsilon-1),\\
 N&=180W+60\epsilon u+\epsilon(3+7\epsilon).
\end{align*}
Multiplying \(\bar w+\delta/2+h B\) by \(|w+\epsilon/6|^2\),
its imaginary part divided by \(\Im w\) is
\(-(\ell\epsilon D+\delta N)/(144\delta)\). Therefore
\(\ell\epsilon D+\delta N=0\). Completion of squares gives
\begin{align*}
 N&=180|w+\epsilon/6|^2+\epsilon(3+2\epsilon)>0,\\
 \delta N-\epsilon D
 &=60(2\epsilon+3)|w+\epsilon/5|^2
    +\epsilon(6\epsilon^2/5+19\epsilon/5+3)>0 .
\end{align*}
Since \(\ell>-1\), these inequalities force \(D<0\) and
\(\ell=-\delta N/(\epsilon D)>0\).
After this substitution,
\[
 \Re(\bar w+\delta/2+hB)=-\frac{P_0}{8\epsilon D},
\]
where
\begin{align*}
P_0={}&720W^2-72W\epsilon^2+480W\epsilon u-168W\epsilon
       -3\epsilon^4-24\epsilon^3u-2\epsilon^3\\
     &+48\epsilon^2u^2-24\epsilon^2u+5\epsilon^2
       -48\epsilon u^2+48\epsilon u.
\end{align*}
These two real and imaginary identities follow by expansion of
\eqref{eq:mid-q-3}. Equation~\eqref{eq:mid-q-4} now gives
\(X=8N/P_0\), so \(P_0>0\).

We claim \(8N-25P_0>0\) on the required range.
The inequality \(D<0\) gives
\[
 u^2\leq W<W_{\max}=(1-\epsilon)u/5+\epsilon(1-\epsilon)/60.
\]
In particular \(u>-\epsilon/12\), since \(W_{\max}>0\).
Also \(u<1/5\): the quadratic
\(60u^2-12(1-\epsilon)u-\epsilon(1-\epsilon)\) is positive
at \(1/5\) and increasing thereafter.
The polynomial \(8N-25P_0\) is concave in \(W\), so it
suffices to check \(W=u^2\) and \(W=W_{\max}\).
At \(W_{\max}\) it is
\[
 4(\epsilon+6u)(2\epsilon+3)
       (5\epsilon^2+10\epsilon u-5\epsilon-10u+4)>0,
\]
since \(\epsilon+6u>\epsilon/2\) and the last factor is at
least \(4-5/4-2>0\).

For \(W=u^2\), write \(8N-25P_0\) as
\[
 (1440+5400\epsilon+600\epsilon^2)u^2
 -12000\epsilon u^3-18000u^4+L u+J,
\]
where
\(L=120\epsilon(5\epsilon^2+5\epsilon-6)\) and
\(J=\epsilon(75\epsilon^3+50\epsilon^2-69\epsilon+24)\).
For \(|u|\le1/10\) this is at least \(A u^2+L u+J\) with
\(A=1260+4200\epsilon+600\epsilon^2\). Its discriminant is
negative, since
\[
 (4AJ-L^2)/(240\epsilon)
 =-750\epsilon^5+2750\epsilon^4+6485\epsilon^3
       +60\epsilon^2-1929\epsilon+504>0.
\]
For \(1/10\le u\le1/5\) use
\(A=720+3000\epsilon+600\epsilon^2\).
Then \(A u^2+L u+J\) is increasing on this interval, since
\(A/5+L=144-120\epsilon+720\epsilon^2+600\epsilon^3>0\),
and its value at \(1/10\) is
\[
 75\epsilon^4+110\epsilon^3-3\epsilon^2-18\epsilon+36/5>0.
\]
For verification of the two displayed polynomial positivities,
in each polynomial substitute \(\epsilon=(l+h t)/(1+t)\)
and clear \((1+t)^{\deg}\) and positive rational denominators.
The descending coefficient lists on the three subintervals
are as follows:
\[
\begin{array}{c|r}
(l,h)&\text{coefficients for the degree-five polynomial}\\ \hline
(0,1/4)&70061,417612,1156648,1594752,1043328,258048\\
(1/4,1/2)&262944,785392,927328,618192,282998,70061\\
(1/2,1)&113920,370432,456588,266402,74168,8217\\[2pt]
(l,h)&\text{coefficients for the degree-four polynomial}\\ \hline
(0,1/4)&5791,21304,37776,31104,9216\\
(1/4,1/2)&20336,49664,48936,25024,5791\\
(1/2,1)&13696,30544,24996,9044,1271
\end{array}
\]
This proves the claim. It gives \(X=8N/P_0>25\).
The coefficient \(x\) is, up to sign, the sum of the
five reciprocal roots of \(P\), of modulus strictly below one.
Consequently \(X<25\), a contradiction. Hence \(w\) is real.
If \(w=-\epsilon/6\), equation \eqref{eq:mid-q-3} would force
\(\ell=-3\delta/\epsilon<-1\). Otherwise that equation forces
\(h\) real. Then \(d\), \(e\), and \(c\) are real as well.
\end{proof}
\begin{proposition}\label{prop:middle-positive-phase}
For \(2<a<3\), every extremizer for \(C(5,p)\) has a polynomial
factor with strictly positive coefficients after rotations.
\end{proposition}

\begin{proof}
Use the notation and first four jets in Lemma~\ref{lem:middle-real-phase}.
First \(x\ne0\) before rotation: otherwise jets one, four, and
two give \(e=b=c=0\), respectively, making \(P\) a polynomial
in \(z^5\). Its sharp coefficient could then be at most
\(C(1,p)<C(5,p)\), where strict monotonicity follows from the
outer theorem by multiplying a first-order maximizer by \(z^4\).
Rotate so that \(x>0\). All remaining coefficients are real by
Lemma~\ref{lem:middle-real-phase}, and \(h=x/(ad)\).

Set \(H=\epsilon(1-\epsilon)/60>0\),
\(m=\epsilon/3\), \(J(w)=w^2+mw-H\), and \(A_2=\binom a2\).
The fifth Taylor coefficient of \(P^{a+1}\), divided by \(a+1\), is
\[
 d+axe+(ab+A_2 x^2)c+A_2x^5 J(w).
\]
The first three terms are aligned with \(d\) by the first two
jets. For changes to the phase of \(b\), write \(u_1=|b|/x^2\)
and \(Y=A_2x^5(u_1^2e^{2i\theta}+mu_1e^{i\theta}-H)\).
Rotating \(d,e,c\) together forces \(Y\) to have the phase of
\(d\) whenever \(Y\ne0\).
In fact \(Y\ne0\) at the extremum. If \(J(w)=0\), then
\(Y'(\theta)=i A_2x^5(2w^2+mw)\ne0\), since
\(J(0)<0\) and \(J(-m/2)<0\).
At the real coefficient point the mixed second derivative of
the squared modulus along the joint phase of \(d,e,c\) and
the phase of \(b\) is nonzero, while its diagonal entry in
the joint-phase direction is zero. This contradicts negative
semidefiniteness. Hence \(\operatorname{sgn}(h)=\operatorname{sgn}J(w)\).

When \(w\ne0\), allow the phase of \(b\) to vary, and rephase
\(d,e\) together and \(c\) independently to align all terms
with \(Y\). This can be done continuously near the original
point. If \(ab+A_2x^2=0\), then \(c=0\) and the corresponding
term can be omitted. All these local coefficient variations
are admissible by strict boundary separation. Consequently the
local maximum in \(\theta\) is also a local maximum of
\[
 |d+axe|+|c|\,|ab+A_2x^2|+|Y|.
\]
In \(X_1=\cos\theta\), the middle term is nondecreasing. If
\(w<0\), the original point has \(X_1=-1\) and the right
derivative of \(|Y|\) must be nonpositive. Its sign is the
sign of \(m(w^2-H)-4Hw\). It follows that
\(w^2-(1-\epsilon)w/5-H\le0\), which forces \(w>-\epsilon/12\):
the left side is positive at \(-\epsilon/12\) and decreasing
up to zero on the negative axis. In particular \(w^2<H\) and
\(J(w)<0\). For \(w<0\) we therefore have \(h<0\).
Also the middle term has nonzero factor since
\(w>-\epsilon/12>-\delta/2\).
For \(w>0\) this factor is automatically nonzero.

We now exclude \(h<0\) for either sign of \(w\ne0\).
The preceding bounds ensure \(w+\epsilon/6>0\).
The exact endpoint derivative condition for the full objective
displayed above is
\[
 w\bigl[J(w)/h + m(w^2-H)-4Hw\bigr]\ \geq\ 0.
\]
Indeed the derivative of the middle term at the original
phase is \(|d|aA_2x^4 u_1\); divide the derivative of the
full objective by \(A_2x^5 u_1/|J(w)|\) and use
\(|J(w)|/|h|=J(w)/h\).
Substitute \eqref{eq:mid-q-3} and the definitions of \(D,N\)
from Lemma~\ref{lem:middle-real-phase}, with \(W=w^2\).
Direct expansion gives
\[
 \delta(w+\epsilon/6)J(w)
  +(\ell w-\delta/2)\bigl[m(w^2-H)-4Hw\bigr]
     =\frac{w}{180}(\ell\epsilon D+\delta N).
\]
Since \(h<0\), the denominator \(\ell w-\delta/2\) is negative.
Thus \(\ell\epsilon D+\delta N\le0\).
The inequalities \(N>0\), \(\delta N-\epsilon D>0\) and
\(\ell>-1\) from the preceding lemma now force
\(D<0\) and \(\ell\ge\ell_0=-\delta N/(\epsilon D)>0\).

The same bound holds when \(w=0\). In that case
\(h=-3/\epsilon<0\) by \eqref{eq:mid-q-3}.
To check the local maximum condition, set \(w=i\eta\)
for real small \(\eta\) and rephase as above, keeping all
other moduli fixed. Set \(U=\|P\|_2^2\) before the variation.
The norm squared of the varied polynomial is \(U+x^4\eta^2\).
The unnormalized coefficient modulus divided by \(a+1\)
at the original point is \(U/((a+1)|d|)\), by the stationary
value formula. Its increment to second order in \(\eta\),
before dividing by the power of the norm, is
\[
 \left(\frac{|d|a^2x^4}{2}
       +A_2x^5\left(1+\frac{m^2}{2H}\right)\right)\eta^2.
\]
Here the first term comes from
\(|c||ab+A_2x^2|\), and the second from
\(|J(i\eta)|=H+(1+m^2/(2H))\eta^2+O(\eta^4)\).
The normalization subtracts \(x^4/(2|d|)\) from this
quadratic coefficient, up to an overall positive multiplier.
Dividing the result by \(x^4a^2|d|/2\) shows that a maximum
requires \(-\ell+\delta|h|(1+m^2/(2H))\leq0\).
Since \(|h|=3/\epsilon\), this is precisely
\(\ell\ge-\delta N/(\epsilon D)=\ell_0\) at \(w=0\).

It remains to rule out the bound \(\ell\ge\ell_0\) in these
cases. Define the affine function
\[
 Q(\ell)=w+\delta/2+
       \frac{(\ell w-\delta/2)B}{\delta(w+\epsilon/6)}
\]
so \eqref{eq:mid-q-4} gives \(X=\ell/(\delta Q(\ell))\).
We have \(Q(0)>0\): for \(w\ge0\) the value of
\(2(w+\epsilon/6)Q(0)\) is
\((3\epsilon^2+20\epsilon w+5\epsilon+36w^2+24w)/24>0\).
For \(-\epsilon/12<w<0\) it follows from
\(w+\delta/2>0\) and \(B<0\).
Also \(Q(\ell)>0\) since \(X>0\) and \(\ell\ge\ell_0>0\).
By affinity \(Q\) is positive on \([0,\ell]\), and
\(\lambda/Q(\lambda)\) is strictly increasing there.
At \(\lambda=\ell_0\), the algebra from the preceding
lemma gives \(Q(\ell_0)=-P_0/(8\epsilon D)\), so \(P_0>0\).
Consequently \(X\ge8N/P_0>25\) by the proved inequality
\(8N-25P_0>0\) on \(D<0\). This contradicts \(|x|<5\).

Thus \(w>0\) and \(h>0\), so \(d,b>0\).
The first two jets give \(e,c>0\), as claimed.
\end{proof}


\subsection{The integer anchors and the range above three}

We include the argument for the remaining exponents. When \(a=2\)
or \(a=3\), the power \(s=a+1\) is an integer. Taking moduli of
the coefficients of the normalized factor before taking its
power preserves admissibility and can only increase the fifth
coefficient modulus. At a nonnegative maximum, let \(S\) be
the support of its constant-term-one factor \(P\) and \(J\)
the support through degree five of \(P^a\). The positive
linear term in the expansion gives \(S\subseteq J\), and the
jet gives \(J=5-S\). Hence \(S=J=5-S\). The quadratic term
gives closure of \(S\) under sums at most five, and reflection
gives closure under nonnegative differences. Euclidean division
now gives \(S=\{0,h,\ldots,5\}\) for a divisor \(h\) of five.
The case \(h=5\) was already excluded by strict monotonicity.
Thus every modulus is positive. Equality of the monomial
phases in the triangle bound proves the same conclusion after
rotation for the original complex factor. In detail, for each
\(2\le j\le5\), compare the phases of
\(q_jq_{5-j}q_0^{s-2}\) and \(q_1q_{j-1}q_{5-j}q_0^{s-3}\).
This gives \(\arg(q_j)=j\arg(q_1)\) when \(q_0>0\).

For \(a>3\), combine the coefficient grouping in
Section~\ref{sec:low-hessian} with local phase variations.
At a maximum \(Y\ne0\), and \(d\) and \(Y\) have the same phase.
If \(b=0\), the third jet gives \(ad/x=-(a-2)/3<0\), whereas
\(Y=\gamma x^5>0\). Thus \(b\ne0\).
For \(b=|b|e^{i\theta}\), write
\[
 Y=L_0e^{2i\theta}+M_0e^{i\theta}+H_0,\qquad
 L_0=\alpha x|b|^2,\quad M_0=\beta x^3|b|,\quad
 H_0=\gamma x^5.
\]
The three constants are positive, and
\(M_0^2/(4L_0H_0)=5(a-2)/(3(a-3))>1\).
Therefore the derivative of \(|Y|^2\) with respect to
\(u=\cos\theta\), namely \(8L_0H_0u+2M_0(L_0+H_0)\), is
strictly positive on \([-1,1]\).
At fixed coefficient moduli, rephase \(d,e\) together and \(c\)
independently to align their contributions with \(Y\).
If the coefficient of \(c\) vanishes at the original point,
then \(c=0\) by the second jet and can be held at zero.
The aligned objective is
\(\lvert d+axe\rvert+|c|\lvert ab+\alpha x^2\rvert+|Y|\).
Its middle term is nondecreasing in \(u\), and the last is
strictly increasing. Any phase with \(u<1\) therefore admits a
local improvement. Consequently \(b>0\), and the alignment
and first two jets give \(d,e,c>0\).
This completes the positivity proof for every \(a>1\).

\section{The distinguished simple root}
\label{sec:branch}

The formula simplifies to a single root. Use \(L=2w\) and
\(t=1/y=x/(ad)\). With \(B=a-2\) and \(\delta=a-1\), put
\begin{align*}
 D&=\delta+L+\delta t(L+B/3),&
 A&=\delta(L+\delta)/2,\\
 J&=\delta\bigl(3L^2+6BL+B(B-1)\bigr)/24,&
 U_4&=\delta(L+\delta)(L+B)/4,\\
 U_5&=\delta B\bigl(15L^2+10(B-1)L+(B-1)(B-2)\bigr)/120 .
\end{align*}
The three real jet equations from Theorem~\ref{thm:main} give
\(X=Lt^2/D\). Substituting this into the fourth and fifth
equations gives
\begin{align*}
 F(a,L,t)&=D(D-L)-L^2(At^2+Jt^3)=0,\\
 G(a,L,t)&=LD^2+a\delta L^2t^2D+
                 aL^3(U_4t^4+U_5t^5)-aD^3=0.
\end{align*}
The calculation is reversible when \(L,t,D>0\).
Define
\begin{equation}\label{eq:branch-R}
 \mathcal P(a,L)
    =\frac{\operatorname{Res}_t(F,G)}
            {L^9(L+\delta)^5\delta^8}.
\end{equation}
The denominator divides the numerator as a polynomial.
Multiply by a nonzero constant to give \(\mathcal P\)
leading \(L\)-coefficient \(-273375\).
From a positive pair \(L,t\) with \(D>0\), write
\(v=\sqrt{L/D}\) and
\[
 x=vt,\qquad b=Lx^2/2,\qquad
 c=v(b+\delta x^2/2),\qquad e=vx,\qquad d=v/a.
\]
\begin{theorem}\label{thm:fifth-unique-branch}
For every \(a>1\), let \(L(a)\) be the unique root of
\(\mathcal P(a,L)\) in \(1<L<2\). Let \(t(a)\) be the unique common
root of \(F(a,L(a),t)\), \(G(a,L(a),t)\).
It is positive and gives the unique normalized maximizing factor.
With the coefficient reconstruction already displayed,
\(C(5,p)=q_0^{a-1}/d\). The factor and sharp value are
real analytic throughout \(a>1\).
\end{theorem}

\begin{proof}
The exact certificate below checks three facts by rational
polynomial arithmetic. First, \(\mathcal P\) has degree fifteen in \(L\)
and constant leading coefficient. Under
\(L=(1+2z)/(1+z)\), multiplication by \((1+z)^{15}\)
gives coefficients of \(z^0,\ldots,z^{10}\) strictly negative
for \(1\le a\le3\) and coefficients of \(z^{12},\ldots,z^{15}\)
strictly positive. The coefficient of \(z^{11}\) is unrestricted.
This is certified in the Bernstein basis on \(1\le a\le3\).
For \(a\ge3\), the signs of coefficients \(z^0,\ldots,z^9\)
and \(z^{12},\ldots,z^{15}\) stay fixed. The coefficients of
\(z^{10}\) and \(z^{11}\) each change from negative to positive
at a single real root, respectively in \((16,17)\) and \((6,7)\),
as certified by rational Sturm counts. Thus for each \(a>1\)
there is exactly one sign change. Descartes' rule shows
\(L(a)\) exists and is simple.

Second, the discriminant of \(\mathcal P\) in \(L\) has exactly sixteen
distinct zeros in \(a>1\). Three are \(3/2,2,3\), and the
others lie respectively in \((k,k+1)/10^6\) for
\[
 \begin{gathered}
 1157449,\ 1438000,\ 1863796,\ 2016999,\ 2022448,\ 2030525,\ 2216434,\\
 2789473,\ 3036377,\ 4594794,\ 7212547,\ 28835694,\ 132692553 .
 \end{gathered}
\]
Third, at the rational samples
\[
 \begin{gathered}
 1.10,\ 1.30,\ 1.45,\ 1.70,\ 1.90,\ 2.01,\ 2.02,\ 2.025,\ 2.10,\\
 2.50,\ 2.90,\ 3.01,\ 4.1,\ 5,\ 10,\ 30,\ 150
 \end{gathered}
\]
there is exactly one positive \(L\) passing all the necessary
tests \(t>0,\ D>0,\ x^2(1-L)<5\), where \(t\) comes from the
degree-one subresultant and \(x=t\sqrt{L/D}\).
This survivor belongs to \(1<L<2\). The last test is necessary
because \(x^2(1-L)=x^2-2b\) is the sum of the squares of the
five reciprocal roots of \(P\), all of modulus less than one.
An extremizer exists for each sample and is positive by the
preceding proof, so the single surviving factor is stable.

We explain why these exact samples suffice on the complementary
parameter intervals. A stable positive solution has \(L>0\),
hence at a simple root of \(\mathcal P\) the resultant itself has a
simple zero as a function of \(L\). The Sylvester matrix of
\(F,G\), with their generic degrees \(3,5\) in \(t\), has
rank deficiency one at this point: if its deficiency were
larger, every determinant derivative would vanish. A finite
common root \(t\) gives a null vector
\((1,t,\ldots,t^7)\) in the matrix arranged with the
coefficient rows in ascending order, so there can be only
one such root. The real analytic continuation of \(L\) preserves
this property locally. Indeed, the rank stays constant on
the branch, its kernel is locally analytic, and its first
coordinate stays nonzero after normalization. The homogeneous
forms have a common projective root by the resultant criterion.
A common root at infinity would give the kernel vector with
first coordinate zero. Thus a finite common root persists
and is recovered analytically as the second coordinate of
the normalized kernel. Positivity and stability are open.

Limits of stable positive solutions at an interior parameter
remain stable and positive. To see this, the reciprocal roots
in the unit disk bound the coefficients of \(P\) by binomial
coefficients, and the fifth jet \(d[z^5]P^a=1\) keeps \(d\)
bounded away from zero. Any limiting \(P\) is zero-free in the
open disk and retains its jet. The reflected boundary
identity excludes roots on the unit circle for such a jet:
its proof uses the power \(P^a\) only as a Taylor jet at zero
and \(h\) only as the factored trigonometric polynomial.
If \(x\) tended to zero, so would \(e,b,c\) by the first,
fourth, and second jets. The sparse limit has \(v^2=a\) with
\(v=ad\). In the third equation,
\[
 b(1-v^2-(a-1)vx)
   =\frac{(a-1)vx^2}{2}\bigl(v+(a-2)x/3\bigr),
\]
so \(b=O(x^2)\) along the sequence and \(c=O(x^2)\).
Dividing the fourth jet by \(x\) then gives \(v^2\to1\),
a contradiction. If \(b\) tended to zero with \(x>0\) at
the limit, the third jet gives \(v=(2-a)x/3\). For \(a\ge2\)
this contradicts \(v>0\). For \(1<a<2\), the calculation in
the positivity proof with \(u=0\) excludes this by the fifth
jet and strict root bound alone. The first two jets then
ensure \(c,e>0\) as well. This establishes the limit assertion.

Thus on every open component off the discriminant zeros
the finite cardinality of stable positive stationary factors
is locally constant. Each one continues by its simple
resultant root and no new ones can enter except by convergence
to an existing one, where that continuation is unique.
The exact samples give cardinality one on each component.
There it comes from \(L(a)\) by the same continuation.
At a discriminant zero the distinguished branch still has
simple \(L\). Compactness and the limit argument give a stable
positive factor there. The local Sylvester argument continues
it through the exceptional value. It remains globally maximizing
there by continuity: every fixed stable polynomial is an
admissible test for all nearby exponents, with a continuous
coefficient value.

Finally, a second global maximizer at an exceptional parameter
would give a second isolated maximum in the positive real
coefficient chart. Isolation follows from the finite candidate
list. Choose a small closed ball around it in that chart,
disjoint from the distinguished branch and with a strict
objective gap at the boundary. The gap persists for nearby
parameters by continuity, producing an interior local maximum
in the same ball. The sphere multiplier equation is
\([z^{5-j}]Q^{s-1}=([z^5]Q^s)q_j\) for \(0\le j\le 5\),
where \(Q=P/\|P\|_2=\sum q_jz^j\). It follows by differentiating
the coefficient on the real unit sphere and using homogeneity.
The equation for \(j=5\) and the five other equations give the jet. It is
stable, positive, and distinct from the distinguished branch,
which contradicts cardinality one at the nearby regular
parameters. This proves the claimed uniqueness, including at
the exceptional values. Analyticity follows from the simple
\(L\)-root and Sylvester continuation on all of \(a>1\).
\end{proof}

\section{The fifth coefficient near \texorpdfstring{\(p=1\)}{p=1}}

The next result determines the extremizer on a left neighborhood of
\(p=1\) and identifies its limiting factor. As usual, write
\(a=2/p-1\) and \(\delta=a-1\). The star on a degree-five polynomial
means \(P^*(z)=z^5\overline{P(1/\bar z)}\).

\begin{theorem}\label{thm:fifth-near-one}
There is \(\varepsilon>0\) such that the fifth coefficient extremizer,
for \(1-\varepsilon<p<1\), is unique up to argument rotation and a
unimodular constant. Its normalized factor is \(Q=q_0P\), where
\[
 \begin{split}
 P(z)&=1+xz+bz^2+cz^3+ez^4+dz^5,\\
 q_0&=(1+x^2+b^2+c^2+e^2+d^2)^{-1/2},\\
 d&=\frac{\sqrt{1-\delta h}}{a},\qquad
 c=d\left(ab+\binom a2 x^2\right),\qquad e=axd .
 \end{split}
\]
All its coefficients are positive. The parameters \((x,b,h)\) are
real analytic through \(\delta=0\).
For \(\delta>0\), they are the unique nearby real solution of
\[
 d[z^j]P^a=[z^{5-j}]P\qquad (j=3,4,5),
 \qquad C(5,p)=q_0^\delta/d.
\]
At \(\delta=0\) the solution is specified by
\[
 \begin{gathered}T(x_*)=0,\qquad
 \frac{15737531359}{10^{10}}<x_*<
 \frac{15737531360}{10^{10}},\\
 b_*=\frac{N(x_*)}{D(x_*)},\qquad
 h_*=\frac{x_*^2/2+x_*b_*-x_*^3/6}{b_*}.\end{gathered}
\]
where the three polynomials are given below. In particular,
\[
 \lim_{p\uparrow1}\frac{\log C(5,p)}{1/p-1}
   =2+h_*-\log\!\bigl(2(1+x_*^2+b_*^2)\bigr)
   =1.229188832\ldots .
\]
\end{theorem}

\begin{proof}
Normalize the extremal factor so all coefficients of \(P\) are
positive. Its reflected stationary identity is \(P^a\equiv P^*/\bar d\pmod{z^6}\).
The first two reflected coefficients give \(e,c\) as in the statement.

We first describe the limit along any subsequence \(p\uparrow1\).
Let \(W=|Q|^2\) and \(\operatorname{Ent}(W)=\int W\log W\,dm\)
for normalized arc measure. For any norm-one polynomial \(Q_0\)
with all its roots on the unit circle, \(Q_0^2\) has fifth
coefficient of modulus one. Its normalized \(H^p\) comparison
gives
\[
 \log C(5,p)\ge -\frac1p\log\int |Q_0|^{2p}\,dm
      =(1/p-1)\operatorname{Ent}(|Q_0|^2)+O((1-p)^2).
\]
For the extremal \(Q\), the \(H^1\) coefficient estimate gives
\[
 C(5,p)\le \int |Q|^{2/p}\,dm
      =1+(1/p-1)\operatorname{Ent}(|Q|^2)+O((1-p)^2).
\]
The errors are uniform because the densities are bounded by six.
The constant term of every \(Q\) is at least \(2^{-5}\) by the
reciprocal root bound, so the limit of a convergent subsequence
of factors is nonzero and has no disk zeros. The two inequalities
give \(C(5,p)\to1\). Consequently the limiting factor satisfies
\(|[z^5]Q^2|=1\), and equality in Cauchy--Schwarz gives
self-reciprocity up to a phase. Reciprocity and the absence of
disk zeros put all five roots on the unit circle. The inequalities
then show that the limiting factor maximizes entropy over all
unit-root factors and that the entropy maximum gives the stated
logarithmic limit.

At such a limit \(d\to1\) by positivity and reciprocity, and
\(U=2(1+x^2+b^2)\), using \(x,b\) for their limiting values.
The relation \(C(5,p)=q_0^\delta/|d|\) implies boundedness of
\(h=(1-a^2|d|^2)/\delta\), with limiting value
\(\operatorname{Ent}(|Q|^2)-2+\log U\). Substitute \(e,c\) from
the first two jets into the last three, multiply by \(d\),
divide by \(\delta\), and take the limit. The resulting equations
are
\begin{align}
 hb&=x^2/2+xb-x^3/6,\label{eq:fifth-end-j3}\\
 hx&=xb+b^2/2-x^2b/2+x^4/12,\label{eq:fifth-end-j4}\\
 h+1&=x^2+b^2-(x^2b+xb^2)/2+x^3b/3-x^5/20.\label{eq:fifth-end-j5}
\end{align}
For example, the terms of degrees three through five in
\(P\log P\) give the first-order corrections to \(P^a\).
If \(x=0\), \eqref{eq:fifth-end-j4} forces \(b=0\), and
\eqref{eq:fifth-end-j5} gives \(h=-1\), hence entropy \(1-\log2\).
If \(b=0\) and \(x>0\), \eqref{eq:fifth-end-j3} forces \(x=3\).
Then \eqref{eq:fifth-end-j4} gives \(h=9/4\) while
\eqref{eq:fifth-end-j5} gives \(h+1=9-243/20<0\).
Thus it remains to consider \(x,b>0\). Eliminate \(h\) using \eqref{eq:fifth-end-j3}.
After clearing denominators, the remaining equations become
\begin{align*}
 f_0={}&-6b^3+6b^2x^2-12b^2x-bx^4+12bx^2-2x^4+6x^3=0,\\
 g_0={}&30b^3x-60b^3-20b^2x^3+30b^2x^2+3bx^5\\
       &\hspace{1.5em}{}-60bx^2+60bx+60b-10x^3+30x^2=0 .
\end{align*}
Their resultant in \(b\) is \(144x^3(x-3)T(x)\), and the linear
subresultant is \(24(D(x)b-N(x))\), where
\begin{align*}
 T(x)={}&x^{15}-30x^{14}+300x^{13}-1070x^{12}-1485x^{11}
         +23860x^{10}-17500x^9\\
 &-424200x^8+1451400x^7-905400x^6-2587800x^5
         +3915000x^4-634500x^3\\
 &-1134000x^2+216000x+108000,\\
 D(x)={}&x^{10}-30x^9+615x^8-6600x^7+35340x^6-96660x^5
        +116550x^4\\
 &-21600x^3-48600x^2+10800x+5400,\\
 N(x)={}&70x^{10}-1050x^9+5790x^8-17610x^7+35100x^6-42300x^5
        +20700x^4\\
 &+3600x^3-2700x^2 .
\end{align*}
The gcd at \(x=3\) is proportional to \(b\), so \(T(x)=0\).
Since the five reciprocal roots of the unit-root polynomial
have modulus one, \(0<x\le5\), and \(T(5)\ne0\).
The polynomials \(T,T'\) are coprime. The Sturm variations at \(0\)
and \(5\) are \(10\) and \(4\), giving six roots with the following
intervals and tests. Each interval is \((m,m+1)/10^{10}\).
The lower endpoint signs of \(T\) are \((+,-,+,-,+,-)\), and the
upper endpoint signs are their opposites. The polynomial \(D\)
has no zero on these intervals.
\[
\begin{array}{r|l}
 m & \text{Tests with } b=N(x)/D(x)\\\hline
  5485556808 & b<0\\
  9453541217 & b<0\\
 15737531359 & b>0\\
 17087628460 & b<0\\
 21037801384 & b<0\\
 40430000373 & 5+b-3x<0
\end{array}
\]
These inequalities follow by rational interval Horner evaluation.
The last root is impossible because
\(P/(1+z)=z^2q(z+z^{-1})\) with
\(q(t)=t^2+(x-1)t+b-x-1\). Unit roots require
\(q(-2)=5+b-3x\ge0\).
The unit-root comparison factor with \(x=3/2,b=2\)
has entropy \(9041/2320-\log(29/2)>1\).
For the latter claim, \(\exp(9041/2320-1)>29/2\) already holds
using the first ten terms of the exponential series.
This excludes the entropy \(1-\log 2\) of the \(x=0\) case.
Thus every limit is the single normalized factor determined
by the third root. Its entropy is \(2+h_*-\log U_*\).
For example, the comparison entropy itself follows from
\[
 \operatorname{Ent}(|Q|^2)
 =-\log U+\frac{2}{U}[z^5](P^2\log P),
\]
which holds by self-reciprocity and integration on the circle.
The logarithm at its zeros can be justified by radial limits
of its analytic power series, or by a direct Fourier evaluation
of the finite coefficients against the trigonometric polynomial
\(|P|^2\).

Finally write \(d=\sqrt{1-\delta h}/a\) and eliminate \(c,e\)
using the first two jets. The three remaining real jets,
multiplied by \(d\) and divided by \(\delta\), extend real
analytically through \(\delta=0\), since their unscaled
residuals vanish there identically in \((x,b,h)\).
Their Jacobian at the endpoint is nonsingular. Indeed
\eqref{eq:fifth-end-j3} solves for \(h\) when \(b>0\).
The linear subresultant and resultant, polynomial combinations
of \(f_0,g_0\), give \(D(x)b-N(x)=0\) and the simple root
\(T(x_*)=0\), with \(D(x_*)\ne0\).
These two combinations have independent gradients, so \(f_0,g_0\)
do as well. The analytic implicit function theorem gives a unique
nearby branch. Every global maximizer near \(p=1\)
is on it, since every convergent subsequence at the endpoint
has the same normalized limit, including \(h\).
This proves the theorem.
\end{proof}

\section{Two algebraic degrees}

At the exponents \(p=4/7\) and \(p=4/5\), respectively, the
sharp fifth coefficient has degree \(60\) and \(52\) over
\(\mathbb Q\). Neither value is expressible by radicals.
The approximate values are \(2.3025752118754219893\ldots\) and
\(1.3515024137960917123\ldots\).

To prove the arithmetic assertions, let \(f\) be the primitive
factor of the distinguished \(L\)-polynomial \(\mathcal P(a,L)\)
containing its root in \((1,2)\). At \(a=5/2\), \(f\)
is the primitive specialization of \(\mathcal P\), of degree \(15\).
At \(a=3/2\), divide \(\mathcal P(3/2,L)\) by \((2L-1)^2\) and take
the primitive part, of degree \(13\).
The modular factorizations are
\[
 \begin{array}{c|c|c}
 a & \text{irreducible reduction} &
       \text{square-free factor degrees at a second prime}\\ \hline
 5/2 & 73 & (2,13)\ \text{at }101\\
 3/2 & 23 & (2,11)\ \text{at }29 .
 \end{array}
\]
Thus both Galois groups are transitive and contain a
transposition, by taking the eleventh or thirteenth power of
the corresponding permutation. Transitivity in degree \(13\)
implies primitivity. In degree \(15\) the group also contains a
\(13\)-cycle, incompatible with a block system of blocks of
size three or five. Thus both groups are primitive. Conjugates
of a transposition in a transitive primitive group generate a
symmetric group, since the corresponding invariant graph of
edges must be connected. The Galois groups are \(S_{15}\) and
\(S_{13}\).

Let \(K=\mathbb Q(L)\), \(V=(ad)^2\), and \(U=\|P\|_2^2\).
The linear subresultant expresses \(t\) as a rational function
of \(L\) in \(K\). Reconstruction then expresses \(V,U\) in
\(K\) and gives
\[
 C(5,p)^4=\frac{a^4}{V^2U^{2(a-1)}}.
\]
At the irreducible prime in each table row, this expression is
integral and has nonconstant, nonsquare residue. These facts
are checked by the exact code below. Nonconstancy excludes
rationality, and primitivity of the Galois group implies
\(\mathbb Q(C(5,p)^4)=K\), since the point stabilizer is
maximal. Nonsquareness in the residue field excludes a square
in \(K\). In fact, if a square root had a coefficient on the
power basis with negative valuation at the indicated prime,
multiplying it by the least appropriate prime power would
give nonzero residue with square zero. Otherwise its residue
would square to the displayed nonsquare. Both are impossible.

For a real field \(K\) and positive nonsquare \(A\in K\), the
polynomial \(Z^4-A\) is irreducible over \(K\). A linear factor
would make \(A\) a square. A real quadratic factor containing
the positive real root must also contain its negative,
making the square of that root an element of \(K\). A cubic
factor would leave a linear factor. Applying this with
\(A=C(5,p)^4\) gives degrees \(4\cdot15\) and \(4\cdot13\).
A radical expression would give a radical expression for
\(C(5,p)^4\), whose normal closure has nonsolvable group
\(S_{15}\) or \(S_{13}\).

\section{Scope and comparison with prior work}

The norm in Theorem~\ref{thm:main} is the original unrestricted
Hardy coefficient norm. The new reduction resolves the full exponent
range in degree five by explicit polynomial equations, the selection of a single
simple root, and the exclusion of nonpositive coefficient phases. The earlier general semialgebraic method in \cite{Low}
already gave finite cell-based representations. The present formula
uses a fixed univariate elimination and phase estimates specific
to degree five.

The literature comparison covers the papers of Brevig and Saksman
and of Brevig and S{\o}vik, the earlier manuscripts \cite{Low,High},
the referenced VMRC snapshot, and targeted searches for a sharp
fifth-coefficient formula. No identical full-range formula or
phase-exclusion proof appeared in the sources reviewed. This
comparison is limited to those sources.

\appendix
\section{Symbolic and rational certificates}
\subsection{Degree of the real eliminant}
\begin{verbatim}
import sympy as S

B, w, y = S.symbols("B w y")
a, d = B+2, B+1
L = y*(w+d/2)+d*(w+B/6)
Q = w*w+B*w+B*(B-1)/12
R = B*w*w/2+B*(B-1)*w/6+B*(B-1)*(B-2)/120
H4 = 2*y*L**2-2*y**2*w*L-d*w**2*(2*y*(w+d/2)+Q)
H5 = (a*y*y*L**3-y**3*w*L**2-a*d*y*y*w*w*L
      -a*d*w**3*(y*(w+d/2)*(w+B/2)+R))
res = S.Poly(S.resultant(H4,H5,y),w)
raw, rem = res.div(S.Poly(w**9*(w+d/2)**5,w))
assert rem.is_zero and raw.degree()==15
assert S.expand(raw.LC()+d**8)==0
\end{verbatim}
\subsection{Coefficients of the nonreal quintic}
\begin{verbatim}
import sympy as S

B,r=S.symbols("B r")
delta,a=B+1,B+2
theta,k=(2*B+3)/30,B/6
m,h=r*k+delta/2,(r+3*delta/B)/theta
W=S.factor(r**2-m*h)
X=S.factor(W/(delta**2*(1+r)))
K=S.factor((2*r**2/X-delta**2*r+m**2
   -B*(3-2*B)*r**2/36+2*B*delta*r/3)/W)
t=S.factor((m+h*K)/(2*r))
E=(1/a+X*(delta+(B-2)*r/5)
   +X**2*((delta+(2*B+1)*r/5)*(K-2*k*t+k**2)
          +delta*(3*B+4)*(t-k)/5+delta**2*(3*B+4)/20)
   -1-r)
c5=3*B**3*(B-1)**3*(B+2)*(2*B+1)
c4=-60*B**3*(B-1)*(B+1)*(B+2)*(B**2+7*B+2)
c3=-10*B**2*(B+1)**2*(B+2)*(8*B**4+70*B**3+77*B**2-456*B-99)
c2=-20*B*(B+1)**3*(12*B**5+50*B**4-307*B**3-1929*B**2-2556*B-270)
c1=-5*B*(B+1)**4*(48*B**4-422*B**3-6045*B**2-19611*B-19170)
c0=-20*(B+1)**5*(4*B**4-213*B**3-1791*B**2-4725*B-4050)
N=sum(c*r**j for j,c in enumerate((c0,c1,c2,c3,c4,c5)))
D=20*B**2*(B+1)**4*(B+2)*(2*B+3)**2*(1+r)**2
assert S.factor(E)==S.factor(N/D)
print("The nonreal quintic polynomial identity is exact.")
\end{verbatim}
\subsection{Endpoint root isolation}
\begin{verbatim}
from fractions import Fraction as F
import sympy as S

x,b=S.symbols('x b')
T=x**15-30*x**14+300*x**13-1070*x**12-1485*x**11+23860*x**10 \
  -17500*x**9-424200*x**8+1451400*x**7-905400*x**6-2587800*x**5 \
  +3915000*x**4-634500*x**3-1134000*x**2+216000*x+108000
D=x**10-30*x**9+615*x**8-6600*x**7+35340*x**6-96660*x**5 \
  +116550*x**4-21600*x**3-48600*x**2+10800*x+5400
N=70*x**10-1050*x**9+5790*x**8-17610*x**7+35100*x**6-42300*x**5 \
  +20700*x**4+3600*x**3-2700*x**2
f=-6*b**3+6*b**2*x**2-12*b**2*x-b*x**4+12*b*x**2-2*x**4+6*x**3
g=30*b**3*x-60*b**3-20*b**2*x**3+30*b**2*x**2+3*b*x**5-60*b*x**2 \
  +60*b*x+60*b-10*x**3+30*x**2
sr=S.subresultants(f,g,b)
assert S.expand(sr[-2]-24*(D*b-N))==0
assert S.expand(sr[-1]-144*x**3*(x-3)*T)==0
assert S.Poly(S.gcd(f.subs(x,3),g.subs(x,3)),b).monic().as_expr()==b

class I:
    def __init__(self,a,b=None):
        self.a,self.b=F(a),F(a if b is None else b)
    def __add__(l,r):
        if not isinstance(r,I):r=I(r)
        return I(l.a+r.a,l.b+r.b)
    __radd__=__add__
    def __neg__(s):
        return I(-s.b,-s.a)
    def __sub__(l,r):
        return l+-r
    def __rsub__(r,l):
        return l+-r
    def __mul__(l,r):
        if not isinstance(r,I):r=I(r)
        v=[l.a*r.a,l.a*r.b,l.b*r.a,l.b*r.b]
        return I(min(v),max(v))
    __rmul__=__mul__
    def __truediv__(l,r):
        if not isinstance(r,I):r=I(r)
        assert r.a*r.b>0
        return l*I(1/r.b,1/r.a)
    def __pow__(s,k):
        assert k>=0
        r=I(1)
        for _ in range(k):r=r*s
        return r
def ev(p,i):
    a=I(0)
    for c in S.Poly(p,x).all_coeffs():
        a=a*i+F(c)
    return a
def sign(p,t):
    y=ev(p,I(t))
    assert y.a!=0
    return 1 if y.a>0 else -1
def log_lower(u):
    w=(u-1)/(u+1)
    return 2*sum(w**(2*k+1)/(2*k+1) for k in range(24))
assert S.gcd(T,S.diff(T,x))==1
st=S.sturm(T,x)
def v(t):
    signs=[]
    for p in st:
        r=ev(p,I(t)).a
        if r!=0:signs.append(1 if r>0 else -1)
    return sum(a*b<0 for a,b in zip(signs,signs[1:]))
ks=[5485556808,9453541217,15737531359,17087628460,21037801384,40430000373]
assert v(0)-v(5)==len(ks)
assert sign(T,0) and sign(T,5)
for idx,k in enumerate(ks):
    lo,hi=F(k,10**10),F(k+1,10**10)
    X=I(lo,hi); Y=ev(N,X)/ev(D,X)
    assert sign(T,lo)*sign(T,hi)==-1
    assert Y.a*Y.b>0
    if idx in [0,1,3,4]:assert Y.b<0
    if idx==2:assert Y.a>0
    H=(X**2/2+X*Y-X**3/6)/Y
    U=2*(1+X**2+Y**2)
    if idx<2:
        assert H.b+1<log_lower(U.a)
    if idx>=3:
        assert (5+Y-3*X).b<0
    if idx==2:
        # q(t)=t^2+(X-1)t+Y-X-1 has two simple roots in (-2,2)
        assert (5+Y-3*X).a>0 and (1+Y+X).a>0
        assert (X-1).a>-4 and (X-1).b<4
        assert ((X-1)**2-4*(Y-X-1)).a>0
        K=Y/X**2
        L,M,J=K**2,-K/3,I(F(1,30))
        assert (K+X*(4*L*J+M*(L+J))/(L+M+J)).a>0
    print(idx, float(X.a),[float(Y.a),float(Y.b)], [float(H.a),float(H.b)])
# Entropy competitor trial e^(9041/2320-1) > 29/2.
from math import factorial
sigma=F(9041,2320)-1
assert sum(sigma**i/factorial(i) for i in range(10))>F(29,2)
\end{verbatim}
\subsection{The single maximizing root}
\begin{verbatim}
import sympy as S

a, L, t, z = S.symbols('a L t z')
d = a-1
lam = L/d
D = 1+lam+(L+(a-2)/3)*t
A = d*d*(lam+1)/2
B = d**3*lam**2/8+d**2*(a-2)*lam/4+d*(a-2)*(a-3)/24
U4 = d*d*(lam+1)*(L+a-2)/4
U5 = (d**3*(a-2)*lam**2/8 + d**2*(a-2)*(a-3)*lam/12
      +d*(a-2)*(a-3)*(a-4)/120)
F = S.cancel(d**2*(D*(D-lam)-lam**2*(A*t*t+B*t**3)))
G = S.cancel(a*d**3*(lam*D**2/a+d*lam**2*t*t*D
       +lam**3*(U4*t**4+U5*t**5)-D**3))
R = S.cancel(S.resultant(F,G,t)/(L**9*(L+d)**5*d**8))
R = S.Poly(R,L)
assert R.degree()==15 and R.LC().is_Rational
R = S.Poly(-273375*R.as_expr()/R.LC(),L)

# A single simple root between L=1 and L=2, uniformly in a.
T = S.Poly(sum(R.nth(i)*(1+2*z)**i*(1+z)**(15-i)
               for i in range(16)),z)
for i in list(range(11))+list(range(12,16)):
    P = S.Poly(T.nth(i).subs(a,1+2*z),z)
    m = P.degree()
    s = -1 if i<=10 else 1
    assert all(s*sum(P.nth(k)*S.binomial(j,k)/S.binomial(m,k)
               for k in range(j+1))>0 for j in range(m+1))

# Extend the sign pattern to a>=3 by Sturm root counts.
for i in list(range(10))+list(range(12,16)):
    P = S.Poly(T.nth(i),a)
    assert S.sign(P.eval(3))==(-1 if i<10 else 1)
    assert P.count_roots(3,S.oo)==0
for i,lo,hi in [(10,16,17),(11,6,7)]:
    P = S.Poly(T.nth(i),a)
    assert P.eval(3)<0 and P.eval(hi)>0
    assert P.count_roots(3,S.oo)==P.count_roots(lo,hi)==1

# Isolate all discriminant zeros in a>1.
r = S.Rational
brackets = [(r(1157449,10**6),r(1157450,10**6)),
            (r(1438000,10**6),r(1438001,10**6)),
            (r(3,2),r(3,2)),
            (r(1863796,10**6),r(1863797,10**6)),
            (r(2),r(2)),
            (r(2016999,10**6),r(2017000,10**6)),
            (r(2022448,10**6),r(2022449,10**6)),
            (r(2030525,10**6),r(2030526,10**6)),
            (r(2216434,10**6),r(2216435,10**6)),
            (r(2789473,10**6),r(2789474,10**6)),
            (r(3),r(3)),
            (r(3036377,10**6),r(3036378,10**6)),
            (r(4594794,10**6),r(4594795,10**6)),
            (r(7212547,10**6),r(7212548,10**6)),
            (r(28835694,10**6),r(28835695,10**6)),
            (r(132692553,10**6),r(132692554,10**6))]
disc = S.discriminant(R.as_expr(),L)
critical_count = [0]*len(brackets)
for fac,mul in S.factor_list(disc,a)[1]:
    P = S.Poly(fac,a)
    for (lo,hi),mult in P.intervals(eps=r(1,10**8),inf=1):
        if lo==hi==1:
            continue
        k = [j for j,(b,c) in enumerate(brackets) if b<=lo<=hi<=c]
        assert len(k)==1 and mult==1
        critical_count[k[0]] += 1
assert critical_count==[1]*len(brackets)

samples = list(map(r, ['1.1','1.3','1.45','1.7','1.9','2.01',
                      '2.02','2.025','2.10','2.50','2.90',
                      '3.01','4.1','5','10','30','150']))
bounds = [(r(1),r(1))]+brackets+[(S.oo,S.oo)]
assert all(bounds[k][1]<v<bounds[k+1][0] for k,v in enumerate(samples))

def sign_at_root(pol,Q,lo,hi):
    """Exact sign of rational polynomial at the isolated root of Q."""
    if pol.is_zero:
        return 0
    common = S.Poly(S.gcd(pol,Q),L)
    if lo!=hi and common.count_roots(lo,hi):
        return 0
    while lo!=hi:
        low=high=S.S.Zero
        for c in pol.all_coeffs():
            vals = (lo*low,lo*high,hi*low,hi*high)
            low,high = min(vals)+c,max(vals)+c
        if low>0:
            return 1
        if high<0:
            return -1
        lo,hi = Q.refine_root(lo,hi,steps=5)
    return S.sign(pol.eval(lo))

for v in samples:
    Q = S.Poly(R.as_expr().subs(a,v),L)
    Fv,Gv = S.Poly(F.subs(a,v),t),S.Poly(G.subs(a,v),t)
    assert all(S.gcd(Q,S.Poly(P.LC(),L)).degree()==0 for P in (Fv,Gv))
    H = S.Poly(S.subresultants(Fv.as_expr(),Gv.as_expr(),t)[-2],t)
    assert H.degree()==1
    # t=B0/A0 at each simple common root.
    A0, B0 = H.nth(1), -H.nth(0)
    ttest = S.Poly(A0*B0,L)
    dtest = S.Poly((v-1+L)*A0*A0+(v-1)*(L+(v-2)/3)*A0*B0,L)
    ntest = S.Poly(5*dtest.as_expr()-L*B0*B0*(1-L),L)
    passing = []
    for (lo,hi),m in Q.intervals(eps=r(1,10**8),inf=0):
        assert m==1
        if lo==hi==0:
            continue
        assert lo>0
        if all(sign_at_root(test,Q,lo,hi)>0 for test in [ttest,dtest,ntest]):
            passing.append((lo,hi))
    assert len(passing)==1 and 1<passing[0][0]<=passing[0][1]<2
    print(v,'passed',flush=True)
\end{verbatim}
\subsection{Algebraic degree certificates}
Append the next block to the single-root certificate. It uses
the polynomial objects defined there.
\begin{verbatim}
# This certificate uses F, G, R and a, L, t as defined above.
def extra_arithmetic(aa, p, q, n, exponent, c2, cn):
    f = R.as_expr().subs(a, aa)
    if n == 13:
        f = S.cancel(f/(2*L-1)**2)
    f = S.Poly(f,L).clear_denoms()[1].primitive()[1]
    assert f.degree()==n
    f0 = S.Poly(f,L,modulus=p).monic()
    assert f0.degree()==n and f0.is_irreducible
    f1 = S.Poly(f,L,modulus=q).monic()
    g2 = S.Poly(c2,L,modulus=q)
    gn = S.Poly(cn,L,modulus=q)
    assert f1==g2*gn and g2.is_irreducible and gn.is_irreducible
    sr = S.Poly(S.subresultants(F.subs(a,aa),G.subs(a,aa),t)[-2],t)
    assert sr.degree()==1
    def fp(expr):
        terms = S.Poly(expr,L).terms()
        return S.Poly(sum(int(c.p)*pow(int(c.q),-1,p)*L**k[0]
                          for k,c in terms),L,modulus=p)
    def red(f):
        return f.rem(f0)
    def inv(f):
        return red(f).invert(f0)
    one = fp(1)
    l = fp(L)
    tt = red(-fp(sr.nth(0))*inv(fp(sr.nth(1))))
    dd = red(l+fp(aa-1)+fp(aa-1)*tt*(l+fp((aa-2)/3)))
    vv = red(l*inv(dd))
    xx = red(vv*tt*tt)
    b = red(l*xx*fp(S.Rational(1,2)))
    u = red(one+xx+b*b+vv*(b+fp((aa-1)/2)*xx)**2
            +vv*xx+vv*fp(1/aa**2))
    c4 = red(fp(aa**4)*inv(vv**2*u**exponent))
    assert c4.degree()>0
    result, base, k = one, c4, (p**n-1)//2
    while k:
        if k%2:
            result = red(result*base)
        base = red(base*base)
        k //= 2
    assert result == -one

extra_arithmetic(S.Rational(5,2),73,101,15,3,L**2+17*L-20,
    L**13-25*L**12-31*L**11+37*L**10+28*L**9-31*L**8
    +49*L**7-39*L**6-47*L**5+9*L**4+5*L**3+5*L**2+L+47)
extra_arithmetic(S.Rational(3,2),23,29,13,1,L**2+9*L+2,
    L**11+11*L**10-12*L**9+14*L**8-14*L**7+8*L**6
    -9*L**4-8*L**3-6*L+3)
\end{verbatim}

\enlargethispage{\baselineskip}
\small
\begin{thebibliography}{9}
\setlength{\itemsep}{2pt}
\bibitem{BrevigSaksman}
O.~F.~Brevig and E.~Saksman, Coefficient estimates for \(H^p\)
spaces with \(0<p<1\), \emph{Proc. Amer. Math. Soc.} \textbf{148}
(2020), 3911--3924.
\href{https://doi.org/10.1090/proc/14995}{DOI: 10.1090/proc/14995}.

\bibitem{BrevigSovik}
O.~F.~Brevig and \O.~S{\o}vik, Extremals for linear functionals on
\(H^p\) are outer when \(0<p<1\),
\href{https://arxiv.org/abs/2610.21788}{arXiv:2610.21788}, 2026.

\bibitem{Low}
\emph{Sharp third and fourth Taylor coefficients in Hardy spaces
below \(H^1\)}, preceding manuscript, November 2026.

\bibitem{High}
\emph{Sharp fifth and sixth Taylor coefficients in Hardy spaces
below \(H^1\)}, preceding manuscript, November 2026.

\bibitem{VMRC}
Verifiable Mathematical Reasoning Corpus, manuscripts [248] and
[249] on sharp zero-free Hardy coefficients and the Krzy\.z problem,
\href{https://github.com/QEEDDD/Verifiable-Mathematical-Reasoning-Corpus/tree/e73bae9fd0c319129b19aacc5c317f7abf2865ef}{snapshot e73bae9f},
2026.
\end{thebibliography}
\end{document}
